% Official contribution: https://conferences.iaea.org/event/498/contributions/43619/
% Scheduled date: 2026-10-28.
\paragraph{Abstract}
% No substantive abstract published at the time of preparation.
\mbox{}

Georg Schnabel briefly presented the IAEA Fusion Evaluated Nuclear Data Library (FENDL)~\cite{Schnabel2024FENDL} as an example to discuss challenges associated with the maintenance and development of large nuclear data library projects with a long history and many contributors. Two key aspects are traceability of changes over time and the reproducibility of derived files, such as specific files for transport codes, e.g. MCNP. Key information technologies adopted by the FENDL project were briefly discussed: Git and GitHub for traceability, git-annex for dealing with many large files, Apptainer for perfect reproducibility of NJOY2016~\cite{NJOY} processing, and the JADE tool for library verification and validation, developed by Fusion For Energy (F4E), United Kingdom Energy Authority (UKAEA) and partners.

Traceability is also improved by translating legacy nuclear data formats to commonly used Internet container and exchange formats, such as the Extensible Markup Language (XML) and JavaScript Object Notation (JSON). Modern tools and programming languages seamlessly support parsing and serializing these formats, which simplifies workflows, the development of new tools and the integration into agentic workflows. At the community level, the adoption of modern formats lowers technological entrance barriers to newcomers and consequently promotes robust and sustainable inter-generational knowledge transfer.

With regard to nuclear reaction data, the General Nuclear Database Structure (GNDS) format~\cite{GNDS2023} is an important modernization project that currently uses XML as a container format. While GNDS will gradually replace the Evaluated Nuclear Data File (ENDF-6) format~\cite{ENDF6FormatsManual2023} for the exchange and dissemination of evaluated nuclear reaction data, ENDF-6 is still widely used by contemporary library projects as well as their previous versions over the last several decades. Seamless access to these data is therefore also paramount to ensure that knowledge from the past remains accessible and can inform future developments. Removing the technological barrier imposed by the ENDF-6 format was the primary motivation for the development of the Python packages endf-parserpy~\cite{endf-parserpy} and endf-userpy~\cite{endf-userpy} at the IAEA Nuclear Data Section so that traceable and reproducible workflows can also integrate evaluated data provided in the ENDF-6 format.

The endf-parserpy package enables low-level parsing of ENDF data and their conversion to nested Python dictionaries, which are first-class citizens in Python and hence powerful functionality to traverse, analyze and modify them is available out-of-the-box. For example, a nested dictionary with ENDF data can be converted to a JSON file in only a couple lines of code, which enables seamless access to ENDF data from many other programming languages and web applications. Furthermore, Python is one of the most prominent languages for Machine Learning (ML) and data science work, which simplifies the integration of ENDF data into data pipelines and ML applications.
As library maintenance also entails fixing and updating evaluated nuclear data occasionally, endf-parserpy also allows for the serialization of Python dictionaries with ENDF data of appropriate structure to ENDF files. By extension, JSON files with ENDF data can be seamlessly serialized to the ENDF format as well. The Python package is available under the permissive MIT license on GitHub at \url{https://github.com/iaea-nds/endf-parserpy} and comprehensive documentation with many tutorials is provided at \url{https://endf-parserpy.readthedocs.io/en/latest/}.

The other Python package, endf-userpy, provides a simple user-interface with a very limited set of high-level functions that return evaluated nuclear data in a form that can be directly compared with experimental data (or is very close to this goal). Examples of quantities of interest are cross sections, angular distributions and double-differential cross sections for both inclusive reactions, e.g. (n,tot) and exclusive ones, e.g. (n,2n). Under the hood, it traverses nested Python dictionaries as produced by endf-parserpy and assembles and interpolates data stored therein as needed to correctly reconstruct the physical quantity queried by the user. An early version for testing (in alpha status) is available at \url{https://github.com/iaea-nds/endf-userpy}. 


\paragraph{Discussion}
% Source: Webex-

ENDF handling through a formal description of the format, low-level parsing and higher-level access to physical quantities. The recipe-based approach supports nested Python structures and generated C++ parsers; endf-userpy builds on these representations to expose observables. Discussion addressed constructing files from tabulated data, combining information from different sections and remaining gaps in convenience functions and covariance handling. Participants distinguished format checks from physical validation. Issue reports, examples and contributions were discussed as ways to improve the tools and support their maintenance. 
